\documentclass[acmsmall,nonacm,screen,11pt]{acmart}
% Libertine when available (Overleaf); Latin Modern otherwise.
\makeatletter\if@ACM@newfonts\else\usepackage[T1]{fontenc}\usepackage{lmodern}\usepackage{amssymb}\fi\makeatother
\settopmatter{printacmref=false}
\renewcommand\footnotetextcopyrightpermission[1]{}
\usepackage{stmaryrd}
\usepackage{array}
\usepackage{listings}
\usepackage{isabelle-listings}
\usepackage{tikz}
\usetikzlibrary{quantikz2,arrows.meta,positioning,calc}

\DeclareMathOperator{\Tr}{tr}
% quantikz defines \ket, \bra and \braket; we use our own.
\let\ket\relax\let\bra\relax\let\braket\relax\let\ketbra\relax
\newcommand{\ket}[1]{|#1\rangle}
\newcommand{\bra}[1]{\langle #1|}
\newcommand{\ketbra}[2]{|#1\rangle\langle #2|}
\newcommand{\norm}[1]{\|#1\|}
\newcommand{\sem}[1]{\llbracket #1\rrbracket}
\newcommand{\hoaret}[3]{\models_{t}\{#1\}\;#2\;\{#3\}}
\DeclareRobustCommand{\isa}[1]{\textsf{\def\_{\char95\allowbreak}#1}}
\definecolor{backcolour}{rgb}{0.96,0.96,0.93}
\lstdefinestyle{isabellestyle}{
    backgroundcolor=\color{backcolour},
    basicstyle=\ttfamily\footnotesize,
    breakatwhitespace=true, breaklines=true, columns=fullflexible,
    captionpos=b, keepspaces=true,
    numbers=none, showspaces=false, showstringspaces=false, tabsize=2,
    aboveskip=6pt, belowskip=6pt,
    xleftmargin=2pt, xrightmargin=2pt, framexleftmargin=2pt, framexrightmargin=2pt}
\lstset{language={Isabelle}, style=isabellestyle, mathescape=true, escapeinside=~~}
\renewcommand{\lstlistingname}{Listing}
\tikzset{slice/.style={thin, gray!80, dash pattern=on 2pt off 1.5pt, align=center}}

\begin{document}

\title[HHL, Verified to the Gate]{HHL, Verified to the Gate with Quantum Hoare Logic in Isabelle/HOL}
\titlenote{This extended abstract summarizes a 37-page paper in preparation, available from the author.
The Isabelle sources (Isabelle2026 with the AFP) are at \url{https://doi.org/10.5281/zenodo.23812498}.}
\author{Dustin Bryant}
\orcid{0000-0001-9876-3804}
\affiliation{\institution{Independent researcher}\country{USA}}
\email{brydustin@gmail.com}
\authorsaddresses{}

\begin{abstract}
HHL prepares a quantum state proportional to the solution of a linear system. Earlier mechanised
proofs cover its exact case and treat its subroutines as black boxes. We verify HHL in Isabelle/HOL,
from the spectral theorem through quantum Hoare logic down to circuits of one- and two-qubit gates.
Four small rules make the logic work at algorithm scale. In the exact case the repeat-until-success
program is totally correct, and its number of rounds is geometric with mean $1/p$. Beyond the exact
case, the paper's error analysis holds with explicit constants, and for the shifted cycle Laplacian the whole algorithm, oracles included, is a
verified gate list with a closed-form count. Seven statements of the original paper are false. An AI
agent wrote almost all of the 64{,}000 lines, directed by one mathematician.
\end{abstract}

\maketitle

\section{HHL, and Why It Is Hard to Verify}

HHL~\cite{hhl2009} started quantum linear algebra. Given a sparse Hermitian $A$ and a state $\ket b$, it
prepares $\ket x\propto A^{-1}\ket b$ in time polynomial in $\log N$, the condition number $\kappa$ and
$1/\varepsilon$. Phase estimation with $U=e^{iAt_0}$ writes the eigenvalues of $A$ into a clock
register; a rotation controlled by the clock moves the amplitude $C/\lambda$ into an ancilla; the clock
is uncomputed; and the ancilla is measured (Fig.~\ref{fig:hhl}). On outcome~$1$ the system register
holds $A^{-1}b$ up to normalisation. This happens with probability $p=C^2\norm{A^{-1}b}^2$, so the run
is repeated until it succeeds. From a programming-languages perspective, HHL is a probabilistic
program, a loop guarded by a measurement, whose correctness is a statement about the distribution of
its outputs and of its running time.

HHL is hard to verify along every axis at once. Its specification is probabilistic, and the number
of rounds is random. Every step is approximate, since eigenvalues are rarely representable in the clock
and $e^{iAt}$ has to be simulated. Its black boxes, the preparation of $\ket b$ and the controlled
evolution, are where the cost is. The mechanised proofs we know of, in Coq~\cite{zhou2023coqq,wu2024local}, verify
the exact case, in which every eigenvalue lies on the grid of the clock, and take the black boxes as given
unitaries; CoqQ proves partial correctness of the repeat-until-success program~\cite{zhou2023coqq},
following the pen-and-paper proof of Zhou, Yu and Ying~\cite{zhou2019applied}. In Lean,
QEDIA~\cite{ren2026qedia} proves a quantum Hoare logic sound for total correctness, with phase
estimation, repeat-until-success circuits and Hamiltonian simulation as case studies, but not HHL.

\begin{figure}[t]
\centering
\scalebox{0.86}{%
\begin{quantikz}[row sep={0.8cm,between origins}, column sep=0.2cm]
\lstick{clock $\ket{0}^{\otimes k}$} & \qwbundle{k} &  & \gate{H^{\otimes k}}\slice{$\psi_2$}
  & \ctrl{1}\slice{$\psi_3$} & \gate{F^\dagger}\slice{$\psi_4$} & \ctrl{2}\slice{$\psi_5$}
  & \gate{F}\slice{$\psi_6$} & \ctrl{1}\slice{$\psi_7$} & \gate{H^{\otimes k}}\slice{$\psi_8$} & \\
\lstick{system $\ket{0}^{\otimes m}$} & \qwbundle{m} & \gate{U_b}\slice{$\psi_1$} & & \gate{U^{c}}
  & & & & \gate{U^{-c}} & & \rstick{$\propto A^{-1}b$ if $1$}\\
\lstick{ancilla $\ket0$} & & & & & & \gate{R(C/\lambda_c)} & & & & \meter{}
\end{quantikz}}
\\[0.5em]
\renewcommand{\arraystretch}{1.15}\small
\begin{tabular}{@{}l@{\enspace}l@{\qquad}l@{\enspace}l@{}}
$\psi_1$ & $\ket0\otimes b\otimes\ket0$ &
$\psi_5$ & $\sum_j\ket j\otimes\Pi_jb\otimes R(C/\lambda_j)\ket0$\\
$\psi_2$ & $u\otimes b\otimes\ket0$ &
$\psi_6$ & $\sum_j\ket j\otimes N^{-1/2}(U^{j}\otimes I)\,w$\\
$\psi_3$ & $\sum_j\ket j\otimes N^{-1/2}U^{j}b\otimes\ket0$ &
$\psi_7$ & $u\otimes w$\\
$\psi_4$ & $\sum_j\ket j\otimes\Pi_jb\otimes\ket0$ &
$\psi_8$ & $\ket0\otimes w$
\end{tabular}\\[0.2em]
{\small $w=\sqrt{I-C^2A^{-2}}\,b\otimes\ket0+CA^{-1}b\otimes\ket1$}
\caption{The HHL circuit and the state after each of its eight gates, for $N=2^k$, $U=e^{iAt_0}$,
$u=N^{-1/2}\sum_j\ket j$ and every eigenvalue of $A$ on the grid $\lambda_j=2\pi j/(Nt_0)$ (read as
$2\pi(j-N)/(Nt_0)$ for $j\ge N/2$). Gate~3
applies $U^c$ when the clock holds $c$, gate~5 applies $R(C/\lambda_c)$ with
$R(c)\ket0=\sqrt{1-c^2}\ket0+c\ket1$, and $\Pi_j$ projects onto the eigenspace of~$\lambda_j$. The
table is one theorem, \isa{hhl\_steps}, proved by following a single eigenvector rather than expanding
in an eigenbasis.}\label{fig:hhl}
\end{figure}

\paragraph{Contributions} We work in Isabelle/HOL, on QHLProver's formalisation~\cite{liu2019qhl}
of Ying's quantum while-language~\cite{ying2011}, and contribute the following.
\begin{enumerate}
\item \emph{Total correctness.} The repeat-until-success program terminates with probability one in
  the normalised solution, exits after exactly $\ell$ failed rounds with probability $(1-p)^\ell p$, and runs
  $1/p$ rounds on average (Section~\ref{sec:qhl}).
\item \emph{The approximations.} The paper's Theorem~1 holds with the constant
  $c=\pi^2\sqrt{4+2\pi^2}$, and so does its run-time analysis with amplitude amplification.
\item \emph{The subroutines.} Phase estimation, as a matrix identity and as a program; amplitude
  amplification; Hamiltonian simulation by Trotter products over sparse oracles; state preparation;
  and the uniformly controlled rotation; the Fourier transform is the AFP's~\cite{manrique2025qft}.
\item \emph{The gates.} For the shifted cycle Laplacian, everything from the preparation of $\ket b$
  to the flag is a list of one- and two-qubit gates with a closed-form count
  (Section~\ref{sec:gates}).
\item \emph{The errata.} Seven statements of the HHL paper are false (Section~\ref{sec:errata}).
\item \emph{The process.} An AI agent wrote almost all of the 130 theories (64{,}010 lines, no \texttt{sorry}), directed
  by one person (Section~\ref{sec:agent}).
\end{enumerate}

\section{Quantum Hoare Logic at Algorithm Scale}\label{sec:qhl}

Programs of the while-language denote completely positive, trace-non-increasing maps $\sem S$ on
partial density operators. A predicate is an operator $0\le P\le I$, and $\hoaret{P}{S}{Q}$ holds if
$\Tr(P\rho)\le\Tr(Q\,\sem S\rho)$ for every state $\rho$; so $\hoaret{p\cdot I}{S}{Q}$ says that $S$
ends in $Q$ with probability at least $p$.

Four rules carry all of our program proofs. Three are derived in Ying's system, namely \emph{initialisation},
\emph{gate lists}, and their composition for the shape \emph{initialise, apply gates, measure} shared by
phase estimation, one run of HHL and all its inexact variants. With $W$ the product of the gates and
$\psi=W\ket0$ the state they prepare,
\begin{gather*}
\{\bra0Q\ket0\cdot I\}\;\mathit{init}\;\{Q\},\qquad\{W^\dagger QW\}\;U_1;\dots;U_m\;\{Q\},\\
\Bigl\{\Bigl(\textstyle\sum_j\langle M_j\psi|\,Q\,|M_j\psi\rangle\Bigr)\cdot I\Bigr\}\;
\mathit{init};\,U_1;\dots;U_m;\,\mathbf{measure}\ M\;\{Q\}.
\end{gather*}
None is deep, but together they turn every probabilistic claim about HHL into one computation. For
$Q=\ketbra tt$ with $\ket t$ in the success branch, the scalar is the fidelity $|\langle t|\psi\rangle|^2$,
so every error bound of the analysis becomes a triple. The fourth rule, for \emph{resetting loops}, is
proved from the semantics, because QHLProver formalises partial correctness only and Ying's
total-correctness rule needs a bound function~\cite{ying2011}.

\paragraph{The HHL loop} The repeat-until-success program
\[
\mathit{hhl\_rus}\;=\;\mathit{init};\,W;\;\mathbf{while}\ \mathit{ancilla}=0\ \mathbf{do}\ (\mathit{init};\,W)
\]
runs the eight gates $W=U_8\cdots U_1$ of Fig.~\ref{fig:hhl} until the ancilla reads~$1$. Because the body begins with $\mathit{init}$, it maps every input $\rho$ to $(\Tr\rho)\,\ketbra\psi\psi$ for
the one state $\psi=W\ket0=\ket0\otimes w$ of Fig.~\ref{fig:hhl}. A round succeeds with probability
$p=\norm{P_1\psi}^2=C^2\norm{A^{-1}b}^2$ and then leaves $\ketbra tt$; otherwise the next round starts
afresh. The loop is the geometric series $\sum_\ell(1-p)^\ell p\,\ketbra tt=\ketbra tt$, and one lemma about
such loops gives termination, the exit distribution and the expected number of rounds. AutoQ checks
two-qubit repeat-until-success circuits automatically and found bugs in published
ones~\cite{chen2026autoq}; here the loop body is a whole algorithm.

\begin{theorem}[exact HHL, repeat until success]
Let $A$ be Hermitian with every eigenvalue on the grid $2\pi j/(Nt_0)$, $-N/2\le j<N/2$, $j\ne0$, of a clock
with $N=2^k$ values, let $U_b\ket0=b$, and let $0<|C|\le2\pi/(Nt_0)$. Then
$\hoaret{I}{\mathit{hhl\_rus}}{\ketbra tt}$ for $\ket t=\ket0\otimes A^{-1}b/\norm{A^{-1}b}\otimes\ket1$,
the loop exits after exactly $\ell$ failed rounds with probability $(1-p)^\ell p$, and it runs $1/p$ rounds on
average.
\end{theorem}

\begin{lstlisting}[float=b,caption={The theorem in Isabelle, with its context. The registers are clock, system and ancilla. \isa{spec} puts the spectrum on the grid, \isa{Ub} prepares $b$, and \isa{Cb} keeps the rotation angles admissible.},label={lst:rus}]
context
  fixes A Ub :: "complex Matrix.mat" and b :: "complex Matrix.vec" and n
      k :: nat and t0 C :: real
  assumes A: "A ~$\in$~ carrier_mat n n" "hermitian A" and n0: "0 < n" and k:
      "1 ~$\le$~ k" and t0: "0 < t0"
    and spec: "~$\bigwedge$~l. l ~$\in$~ herm_spectrum A ~$\Longrightarrow$~ ~$\exists$~j. 0 < j ~$\wedge$~ j < 2 ^ k ~$\wedge$~ l
        = reg_eig t0 (2 ^ k) j"
    and Ub: "Ub ~$\in$~ carrier_mat n n" "Quantum.unitary Ub" "Ub *~$_{v}$~
        unit_vec n 0 = b"
    and Cb: "~$\bigwedge$~j. j < 2 ^ k ~$\Longrightarrow$~ ~$|$~C / reg_eig t0 (2 ^ k) j~$|$~ ~$\le$~ 1"
begin
theorem hhl_rus_total:
  assumes "C ~$\neq$~ 0"
  shows "state_sig.hoare_total_correct [2 ^ k, n, 2] (1~$_{m}$~ (2 ^ k * (n *
      2))) (hhl_rus A Ub t0 k C) (outer_prod (hhl_target A k b)
      (hhl_target A k b))"
theorem hhl_rus_expected_rounds:
  assumes C0: "C ~$\neq$~ 0"
  defines "p ~$\equiv$~ C~$^{2}$~ * ~$\|$~herm_inv A *~$_{v}$~ b~$\|^{2}$~"
  shows "(~$\lambda$~m. real (Suc m) * (p * (1 - p) ^ m)) sums (1 / p)" and
      "(~$\lambda$~m. (1 - p) ^ m) sums (1 / p)"
\end{lstlisting}

\paragraph{The hypotheses are the specification} Listing~\ref{lst:rus}, with the definitions it
names, is all that a reader has to trust. Measured as in Lean-QuantumAlg~\cite{zhang2026leanquantumalg}, that is 71 lines, 0.35\% of the 20{,}120 lines of our theories that it
imports, with Isabelle/HOL and the AFP audited once (360 lines, 1.06\%, for the gate-level theorem). Writing it down exposed a restriction that informal accounts
hide. The bound on $C$ comes from the program, not from the mathematics. The circuit applies
$R(C/\lambda_j)$ on every clock value, and a gate must be unitary. The bound costs nothing when the
smallest modulus $\lambda_{\min}$ of an eigenvalue is the first grid point $\lambda_1=2\pi/(Nt_0)$, and
otherwise lowers the best possible $p$ by the factor $(\lambda_1/\lambda_{\min})^2$. The programs of the paper's regime clamp the rotation instead, trading
exactness for fidelity triples.

\section{From Black Boxes to Gates}\label{sec:gates}

The programs above take $U_b$, the controlled evolution and the controlled rotation as unitaries. We
refine them until every gate is \emph{elementary}, that is, a one- or two-qubit unitary on adjacent wires.
Every program theorem is stated for all gate lists with a given product, so refinement never reopens
it. For $s$-sparse
$A$ with entries bounded by $L$, whose graph is coloured into $2s-1$ matchings given by oracles, $r$
Trotter steps are within $8(sLt)^2/r$ of $e^{iAt}$ and make $r(6s-2)$ oracle queries.

For the Laplacian of the cycle on $2^m$ vertices, shifted by $\mu\ge0$, the oracles themselves are
compiled, as reversible arithmetic routed to adjacent wires, and so is the preparation of a sparse
nonnegative $b$. One theorem collects the result. It gives an explicit list of elementary gates on
$K=m+k+1+a$ qubits whose output, post-selected on the flag, is within $2\delta/(C\norm x)$ of $x/\norm x$,
where $x$ solves the system for the right-hand side that the gates actually prepare. The error budget
$\delta=\epsilon+2(2^k-1)\eta$ adds the error $\epsilon$ of phase estimation to an error $\eta$ for each
simulated evolution, and the length is an explicit polynomial, $O(4^k(2+\mu)^2t_0\,m^2K/C)$, which is quadratic
in the clock size because first-order Trotter steps must shrink as the clock grows. Black boxes hide the cost. The
algorithm and its programs take 19 theories, and the oracles of this one family take 39 more. Like the
end-to-end proof of
Shor's algorithm in SQIR~\cite{peng2023shor}, the pipeline leaves nothing abstract; unlike verified
oracle compilers~\cite{amy2017reverc,li2022vqo}, it compiles the oracles of one family only.

\section{Seven Errata}\label{sec:errata}

Proving the paper's statements as written, appendix included, refuted seven of them
(Table~\ref{tab:errata}). The most instructive is Lemma~3, the Lipschitz estimate behind
Theorem~1. At $\lambda=1/(2\kappa)$ and $\tilde\lambda=1/\kappa$ its left side is $\frac12$ and its right
side is $\pi^2/32$, for every $\kappa$. The constant must be $\pi^2$, not $\pi^2/2$, and $\pi^2$ is sharp.
None of the errors invalidates a main claim. Most sit where a sketch compresses a computation, such as an
endpoint value, an index, or a number of rounds or samples. A reader checks the argument; a proof
assistant checks the arithmetic, where the sketch was wrong.

\begin{table}[t]
\caption{False statements of the HHL paper (arXiv:0811.3171v3~\cite{hhl2008arxiv}), each refuted in
Isabelle.}\label{tab:errata}
\footnotesize
\renewcommand{\arraystretch}{1.1}
\begin{tabular}{@{}l>{\raggedright\arraybackslash}p{4.3cm}>{\raggedright\arraybackslash}p{6.3cm}@{}}
\toprule
Where & Claim & What is true\\
\midrule
Lemma 3 & constant $c=\pi^2/2$ & false for all $\kappa$; $\pi^2$ holds and is sharp\\
Proof of Thm.~4 & $\frac2\delta\log(2n)/\log\log n$ rounds & too few; $\frac2\delta\ln(2n_0)$ suffice\\
(A34) & dilation has $\kappa^{-1}I\le A\le I$ & impossible, as its diagonal is zero\\
(A37) & first qubit of $\vec y$ decides & top half of $\vec y$ is $0$; permuting the dilation fixes it\\
Thm.~6(2) & precision $2^{-n-2}$ & cannot separate consecutive counts\\
App.~A5 & $1/(3\varepsilon^2)$ samples & too few; $3/\varepsilon^2$ suffice\\
\S IV & preconditioning $(AB)c=Bb$ & neither left nor right preconditioning\\
\bottomrule
\end{tabular}
\end{table}

\section{How It Was Built}\label{sec:agent}

Claude Code~\cite{claudecode} wrote the development under the direction of the author, an
experienced Isabelle user who wrote almost none of the proof text. The foundation, Paulsen's lecture
course~\cite{paulsen2011notes}, was written in April and May 2026, and, after a pause of 147 days, HHL in eleven days
of October, with 258 human directions over 84 distinct hours and 9{,}835 agent tool calls. The agent
wrote and repaired the proofs, and resumed after the pause from its own hand-off note.
The human chose what to prove, checked that the statements meant what the paper meant, and asked, on
28 October, to stop treating the black boxes as given. The kernel rejects wrong proofs, so the agent's
failures lay elsewhere, as in a confident ``no oracles'' report from an audit that had skipped the gate
theories, where two theorems depended on \texttt{skip\_proof}. Checks in the build now catch such failures. The kernel guarantees that the
theorems are true; the human guarantees that they are the right theorems.

\section{Discussion and Future Work}

We would value the community's view on four questions. Should closed forms like ours be built into
tools for quantum Hoare logic? Could verified gate counts like ours back the cost models of languages
for quantum linear algebra, such as Cobble~\cite{yuan2026cobble}? What should
``verified to the gate'' mean once the gate set is finite, as in Clifford+T~\cite{ross2016gridsynth}, and the oracles general? And as agents write
more of our proofs, which checks should a verified quantum library require? Our next steps are clamping the rotation, higher-order product
formulas, and a second prover.

\paragraph{Use of AI} AI agents (Claude Code with Claude Opus~4.7, Fable~5.1 and Opus~5.5; Codex with
GPT-5.4) wrote the development under the author's direction and drafted this abstract; the author takes
full responsibility for its content.

\bibliographystyle{ACM-Reference-Format}
\bibliography{refs}

\end{document}
